\scanpage{49}\mspage{48}

\diagram{\begin{minipage}{\linewidth}
\[
\vcenter{\hbox{\begin{tikzcd}[ampersand replacement=\&,column sep=large,row sep=small]
\mathrm{Ssfig}(F') \arrow[r] \& \mathrm{Ssfig}(F^{\unsure{9|\prime|?}}) \& \text{où }F=f(F') \\
G' \arrow[r,maps to] \& f(\unsure{C|G}') \&
\end{tikzcd}}}
\tag*{\makebox[0pt][r]{\makebox[\displaywidth][l]{(1.54)}}}
\]
\end{minipage}}
% slanted
\margin{\unsure{ce qui|qui} implique que \deleted{\unsure{Il est sté}} \deleted{\unsure{que}} $f(\mathcal F')$ \deleted{\unsure{est}}soit une partie fermée de $\mathcal F$ pour $\leq$}
soit \uline{bijective}. Sans supposer $f$ injective, dans tous les cas à ma connaissance, on aura une application canonique (\unsure{croissante})
\diagram{\begin{minipage}{\linewidth}
\[
\vcenter{\hbox{\begin{tikzcd}[ampersand replacement=\&,row sep=large]
{\deleted{\unsure{\psi|q}}}\ \& \widetilde F \arrow[d,"{\unsure{\varphi|\phi}_{f,F'}=\varphi}"] \\
\& \widetilde F'
\end{tikzcd}}}
\tag*{\makebox[0pt][r]{\makebox[\displaywidth][l]{(1.55)}}}
\]
\end{minipage}}
d'ens ordonnés, de telle façon que (1.54) en soit déduite par l'application en sens inverse
\diagram{\begin{minipage}{\linewidth}
\[
\vcenter{\hbox{\begin{tikzcd}[ampersand replacement=\&,column sep=large,row sep=small]
{} \arrow[d,phantom,"{\left\{\vphantom{\dfrac{\displaystyle F}{\displaystyle F}}\right.}" description] \& {\unsure{\mathcal P|\mathrm P}\unsure{f|{}_{f}}(\widetilde F')} \arrow[r,dash,"{\unsure{\shortmid|>}}"{at end,description,inner sep=0pt}] \& \mathcal P f(\widetilde F)\\
{} \& A \arrow[r,maps to] \& \varphi^{-1}(A)\ ,
\end{tikzcd}}}
\tag*{\makebox[0pt][r]{\makebox[\displaywidth][l]{(1.56)}}}
\]
\end{minipage}}
et de plus, la \deleted{topologie de} relation d'ordre sur $\widetilde F$ est image inv. de celle sur $\widetilde F'$.

Dans le cas où $f$ est injective, et où la condition (1.54) est satisfaite, (1.55) sera un \unsure{izo.|iso.} d'ens. ordonnés. Dans le cas général, on \unsure{déduira|désirera} encore que l'application (1\unsure{:|.}54) est \deleted{\unsure{tout}} au moins \uline{surjective}. \deleted{\unsure{Le}} Mais dans des cas impor\unsure{\deleted{tants}|tants} (comme $\mathcal F\to\mathrm{Fig}(\unsure{\mathcal L|L})$), même si $f$ n'est pas injective, on aura (1.54) bijective pour \unsure{tte|toute} $F$.

On notera que dans les cas où (1.54) peut se déduire d'un \deleted{\unsure{ap}} hom d'ens ordonnés (1.55), celui-ci est unique. En d'autres termes
